% © 2025 Ing. David Jaroš — CC BY-NC-ND 4.0
%
% This work is licensed under a Creative Commons Attribution-NonCommercial-NoDerivatives
% 4.0 International License (CC BY-NC-ND 4.0).
%
% License History: Earlier drafts (up to v0.3) were released under CC BY 4.0.
% From v0.4 onward, all material is released under CC BY-NC-ND 4.0 to protect
% the integrity of the theoretical work during ongoing academic development.

\documentclass[11pt,a4paper]{article}
\usepackage{amsmath, amssymb, amsthm}
\usepackage{hyperref}
\usepackage{geometry}
\geometry{margin=2.5cm}

\newtheorem{theorem}{Theorem}
\newtheorem{lemma}{Lemma}
\newtheorem{proposition}{Proposition}
\newtheorem{definition}{Definition}
\newtheorem{remark}{Remark}

% Proof-level macros
\newcommand{\Lzero}{\textbf{[L0]}}
\newcommand{\Lone}{\textbf{[L1]}}
\newcommand{\Ltwo}{\textbf{[L2]}}
\newcommand{\SE}{\textbf{[SE]}}
\newcommand{\OPEN}{\textbf{[OPEN]}}
\newcommand{\PROVED}{\textbf{PROVED}}
\newcommand{\Bbase}{B_{\mathrm{base}}}
\newcommand{\Balpha}{B_{\alpha}}
\newcommand{\Neff}{N_{\mathrm{eff}}}

\title{Maxwell Limit Bridge\\
\large Unified Biquaternion Theory — Canonical Bridges Series}
\author{Ing.\ David Jaroš}
\date{2025}

\begin{document}
\maketitle

\begin{abstract}
This document is a \emph{navigation bridge}: it does not contain new derivations
but traces the derivation chain showing that Maxwell's equations emerge as a
classical, low-energy, flat-spacetime limit of the Unified Biquaternion Theory (UBT).
All referenced results are located in existing repository files; section numbers
and file paths are given explicitly. The Maxwell limit follows directly from the
QED sector of UBT when the phase field $\varphi$ is treated as constant and the
gravitational curvature is neglected. Status labels are defined at
\texttt{canonical/bridges/QED\_limit\_bridge.tex}.
\end{abstract}

\tableofcontents
\newpage

% ------------------------------------------------------------------
\section{Overview of the Derivation Chain}
% ------------------------------------------------------------------

The full derivation chain from UBT to Maxwell's equations proceeds through
four steps:

\begin{enumerate}
  \item \textbf{Biquaternionic field} $\Theta(q,\tau)$ contains the gauge sector
        via its imaginary quaternionic components.
  \item \textbf{Electromagnetic potential} $A_\mu$ is identified with the
        U(1) connection arising from the phase degree of freedom of $\Theta$.
  \item \textbf{Field-strength tensor} $F_{\mu\nu} = \partial_\mu A_\nu - \partial_\nu A_\mu$
        is defined canonically.
  \item \textbf{Maxwell equations} follow from the Euler--Lagrange equations of
        the QED Lagrangian in the limit of constant phase and vanishing curvature.
\end{enumerate}

\noindent Proof-level labels (see \texttt{canonical/bridges/QED\_limit\_bridge.tex}):
\begin{itemize}
  \item \Lzero{} — exact algebraic identity, no free parameters
  \item \Lone{} — proved given verified assumptions
  \item \Ltwo{} — requires additional lemmas not yet proved
  \item \SE{} — semi-empirical
  \item \OPEN{} — open problem
\end{itemize}

% ------------------------------------------------------------------
\section{Step 1: Electromagnetic Potential from U(1) Phase}
% ------------------------------------------------------------------

\noindent\textbf{Status: \PROVED{} \Lone}

\noindent\textbf{Primary source:}
\texttt{canonical/interactions/qed.tex}, \S2;
\texttt{canonical/interactions/sm\_gauge.tex}, \S3.

\medskip

The biquaternion $\Theta(q,\tau) \in \mathbb{H} \otimes \mathbb{C}$ has a scalar
phase degree of freedom $\varphi(x)$ arising from the right U(1) action:
\begin{equation}
  \Theta \;\mapsto\; \Theta\,e^{i\varphi(x)}.
\end{equation}
The U(1)$_Y$ hypercharge group is identified with this right-action phase
(\texttt{canonical/interactions/sm\_gauge.tex}, Theorem G.3, \Lzero).
In the electromagnetic sector the electromagnetic four-potential is:
\begin{equation}\label{eq:potential}
  A_\mu(x) \;:=\; \frac{1}{e}\,\partial_\mu \varphi(x),
\end{equation}
where $e$ is the elementary charge. Under a local U(1) gauge transformation
$\varphi \to \varphi + e\,\alpha(x)$ one recovers:
\begin{equation}
  A_\mu \;\to\; A_\mu + \partial_\mu \alpha(x),
\end{equation}
which is the standard gauge transformation of the electromagnetic potential.

% ------------------------------------------------------------------
\section{Step 2: Field-Strength Tensor}
% ------------------------------------------------------------------

\noindent\textbf{Status: \PROVED{} \Lone}

\noindent\textbf{Primary source:}
\texttt{canonical/interactions/qed.tex}, \S2;
\texttt{consolidation\_project/appendix\_C\_electromagnetism\_gauge\_consolidated.tex}, \S1.

\medskip

The electromagnetic field-strength tensor is defined as the curvature of the
U(1) connection:
\begin{equation}\label{eq:Fmunu}
  F_{\mu\nu} \;=\; \partial_\mu A_\nu - \partial_\nu A_\mu.
\end{equation}
In the biquaternionic formalism the same object appears as the imaginary part
of the field:
\begin{equation}
  \mathcal{F} \;=\; (\mathbf{E} + i\,\mathbf{B})\cdot\boldsymbol{\sigma},
\end{equation}
where $\mathbf{E}$ and $\mathbf{B}$ are the electric and magnetic field vectors
and $\boldsymbol{\sigma}$ are the Pauli matrices
(\texttt{consolidation\_project/appendix\_C\_electromagnetism\_gauge\_consolidated.tex}, \S1).
The standard components are recovered by:
\begin{equation}
  F_{0i} = E_i, \qquad F_{ij} = -\varepsilon_{ijk}B_k.
\end{equation}
The Bianchi identity $\partial_{[\lambda}F_{\mu\nu]} = 0$ follows automatically
from the antisymmetry of $F_{\mu\nu}$ and is the \emph{homogeneous} Maxwell
equation in disguise.

% ------------------------------------------------------------------
\section{Step 3: QED Lagrangian and Euler--Lagrange Equations}
% ------------------------------------------------------------------

\noindent\textbf{Status: \PROVED{} \Lone}

\noindent\textbf{Primary source:}
\texttt{canonical/interactions/qed.tex}, \S1 and \S3.

\medskip

The QED Lagrangian density in the UBT framework is
(\texttt{canonical/interactions/qed.tex}, eq.\ (1)):
\begin{equation}\label{eq:LQED}
  \mathcal{L}_{\mathrm{QED}}
  \;=\;
  \mathrm{Tr}\bigl[(D_\mu\Theta)^\dagger (D^\mu\Theta)\bigr]
  \;-\;
  \frac{1}{4}\,F_{\mu\nu}F^{\mu\nu},
\end{equation}
where $D_\mu = \partial_\mu - ieA_\mu$ is the covariant derivative.
The photon kinetic term $-\tfrac{1}{4}F_{\mu\nu}F^{\mu\nu}$ is invariant under
gauge transformations and produces the dynamics of the electromagnetic field.

Varying \eqref{eq:LQED} with respect to $A_\nu$ gives the Euler--Lagrange equation:
\begin{equation}\label{eq:ELeq}
  \partial_\mu\frac{\partial\mathcal{L}}{\partial(\partial_\mu A_\nu)}
  = \partial_\mu F^{\mu\nu} = j^\nu,
\end{equation}
where $j^\nu$ is the conserved electromagnetic four-current arising from
the matter sector (see below).

% ------------------------------------------------------------------
\section{Step 4: Maxwell Equations as the Low-Energy Limit}
% ------------------------------------------------------------------

\noindent\textbf{Status: \PROVED{} \Lone}

\noindent\textbf{Primary source:}
\texttt{canonical/interactions/qed.tex}, \S3;
\texttt{consolidation\_project/appendix\_C\_electromagnetism\_gauge\_consolidated.tex}, \S2.

\medskip

\begin{theorem}[Maxwell Limit]
In the limit of constant phase ($\partial_\mu\varphi = 0$), flat spacetime
($g_{\mu\nu} = \eta_{\mu\nu}$), and classical matter fields, the equations of
motion derived from $\mathcal{L}_{\mathrm{QED}}$ reduce exactly to Maxwell's
equations.
\end{theorem}

\begin{proof}[Proof sketch]
\emph{Homogeneous equations.}
The Bianchi identity on $F_{\mu\nu}$ (which follows from its definition as a
curl, eq.\ \eqref{eq:Fmunu}) gives:
\begin{equation}\label{eq:MaxwellHom}
  \partial_{[\lambda}F_{\mu\nu]} = 0
  \quad \Longleftrightarrow \quad
  \nabla\cdot\mathbf{B} = 0,\quad
  \nabla\times\mathbf{E} + \frac{\partial\mathbf{B}}{\partial t} = 0.
\end{equation}

\emph{Inhomogeneous equations.}
The Euler--Lagrange equation \eqref{eq:ELeq}, with the matter current
$j^\nu = e\,\bar\psi\gamma^\nu\psi$ from the Dirac sector of UBT, gives:
\begin{equation}\label{eq:MaxwellInhom}
  \partial_\mu F^{\mu\nu} = j^\nu
  \quad \Longleftrightarrow \quad
  \nabla\cdot\mathbf{E} = \rho,\quad
  \nabla\times\mathbf{B} - \frac{\partial\mathbf{E}}{\partial t} = \mathbf{j}.
\end{equation}
Together, \eqref{eq:MaxwellHom} and \eqref{eq:MaxwellInhom} are Maxwell's
equations in Gaussian units ($c=1$). \qed
\end{proof}

\begin{remark}[Curved Spacetime Extension]
In curved spacetime, $\partial_\mu \to \nabla_\mu$ (covariant derivative) and
the flat-metric contractions $F_{\mu\nu}F^{\mu\nu}$ involve $g^{\mu\alpha}g^{\nu\beta}$.
The covariant form of the inhomogeneous equations becomes
$\nabla_\mu F^{\mu\nu} = j^\nu$, which is derived in
\texttt{canonical/interactions/qed.tex}, \S6, and in
\texttt{consolidation\_project/appendix\_C\_electromagnetism\_gauge\_consolidated.tex}, \S3.
\end{remark}

\begin{remark}[Biquaternionic Form]
The biquaternionic representation $\mathcal{F} = (\mathbf{E}+i\mathbf{B})\cdot\boldsymbol{\sigma}$
makes the duality symmetry $\mathbf{E}\to\mathbf{B}$, $\mathbf{B}\to-\mathbf{E}$ manifest
as multiplication by $i$ (complex rotation in the biquaternionic algebra).
This is detailed in
\texttt{consolidation\_project/appendix\_C\_electromagnetism\_gauge\_consolidated.tex}, \S4.
\end{remark}

% ------------------------------------------------------------------
\section{Current Sector and Charge Conservation}
% ------------------------------------------------------------------

\noindent\textbf{Status: \PROVED{} \Lone}

\noindent\textbf{Primary source:}
\texttt{canonical/interactions/qed.tex}, \S4.

\medskip

The conserved Noether current associated with U(1) gauge invariance is:
\begin{equation}
  j^\mu = e\,\mathrm{Tr}\bigl[\Theta^\dagger\,(D^\mu\Theta) - (D^\mu\Theta)^\dagger\,\Theta\bigr]
  \;\xrightarrow{\text{Dirac limit}}\;
  e\,\bar\psi\,\gamma^\mu\psi.
\end{equation}
Charge conservation $\partial_\mu j^\mu = 0$ follows from the Euler--Lagrange
equations and is exact (\Lone).

% ------------------------------------------------------------------
\section{Summary Table}
% ------------------------------------------------------------------

\begin{center}
\begin{tabular}{lll}
\hline
\textbf{Result} & \textbf{Source} & \textbf{Status} \\
\hline
U(1) phase $\to$ gauge potential $A_\mu$ & \texttt{qed.tex} \S2; \texttt{sm\_gauge.tex} \S3 & \PROVED{} \Lone \\
Field tensor $F_{\mu\nu}$ & \texttt{qed.tex} \S2 & \PROVED{} \Lone \\
QED Lagrangian & \texttt{qed.tex} \S1 & \PROVED{} \Lone \\
Homogeneous Maxwell eqs.\ & \texttt{qed.tex} \S3 & \PROVED{} \Lone \\
Inhomogeneous Maxwell eqs.\ & \texttt{qed.tex} \S3 & \PROVED{} \Lone \\
Covariant form (curved spacetime) & \texttt{qed.tex} \S6; \texttt{appendix\_C} \S3 & \PROVED{} \Lone \\
Biquaternionic $\mathcal{F}$ form & \texttt{appendix\_C} \S4 & \PROVED{} \Lone \\
Charge conservation & \texttt{qed.tex} \S4 & \PROVED{} \Lone \\
\hline
\end{tabular}
\end{center}

\begin{remark}[No Open Problems in Maxwell Sector]
Unlike the fine-structure constant (where $\Bbase$ remains open) or the
gauge group chirality (where left-handedness is only motivated), the Maxwell
sector has no outstanding derivation gaps. All four Maxwell equations follow
directly from the QED Lagrangian and standard variational calculus.
\end{remark}

% ------------------------------------------------------------------
\section{Cross-References}
% ------------------------------------------------------------------

\begin{itemize}
  \item \textbf{QED Lagrangian and running coupling}: \texttt{canonical/bridges/QED\_limit\_bridge.tex}
  \item \textbf{U(1) gauge group derivation}: \texttt{canonical/bridges/gauge\_emergence\_bridge.tex}, \S4
  \item \textbf{GR recovery}: \texttt{canonical/bridges/GR\_chain\_bridge.tex}
  \item \textbf{Full EM analysis}: \texttt{consolidation\_project/appendix\_C\_electromagnetism\_gauge\_consolidated.tex}
  \item \textbf{Curved-spacetime Maxwell}: \texttt{canonical/interactions/qed.tex}, \S6
  \item \textbf{Claim evidence}: \texttt{AUDITS/claim\_evidence\_matrix.md}, row ``Maxwell\_limit''
\end{itemize}

\end{document}
